% Submit to an IEEE conference (DSN). IEEEtran.cls ships with Overleaf and with
% TeX Live; it is not vendored here.
%
% Every number in this draft is either traced to a committed result file or wrapped in
% \TODO. Do not remove a \TODO by guessing: the project has already retracted three
% findings that were asserted before they were measured.
\documentclass[conference]{IEEEtran}
\IEEEoverridecommandlockouts

\usepackage{amsmath,amssymb}
\usepackage{booktabs}
\usepackage{graphicx}
\usepackage{xcolor}
\usepackage{url}

\newcommand{\TODO}[1]{\textcolor{red}{\textbf{[TODO: #1]}}}
\newcommand{\G}{\ensuremath{G}}

\begin{document}

\title{What Does Generated Control Code Actually Do?\\
A Storage-Recovery Benchmark with an Exact Optimum}

\author{\IEEEauthorblockN{\TODO{author list}}
\IEEEauthorblockA{R.V.\ College of Engineering, Bengaluru, India}}

\maketitle

\begin{abstract}
Benchmarks for language-model-generated code ask whether the code runs. They cannot ask
how good its behaviour is, because for the open-ended tasks of interest no one knows the
best attainable behaviour. We present a benchmark where that question has an exact
answer. The task is recovery-bandwidth scheduling during a storage rebuild: a closed-loop
control problem over production block traces, with four discrete actions and a state
small enough that the optimal policy is computed exactly by backward induction. Policies
are scored as the fraction of the attainable improvement they capture over a tuned
rule-based controller, so the score is dimensionless and bounded above by a reference
that is optimal rather than merely best-known.

Three results follow. First, control policies written by current models run, read
plausibly, and perform worse than a one-line threshold rule on a single input: across
five hosted frontier models the best median score is $-2.62$ against the rule's $-0.62$,
and four of the five never exceed it on any of eight held-out environments.
Second, iterative refinement, in which a model sees its own policy's measurements and
revises, yields no advantage over blind sampling at a matched evaluation budget: carrying
the full conversation scores $-1.43$ against random search's $-1.445$, while deleting the
model's memory of earlier rounds \emph{improves} it to $-0.43$. Third, inside a population search at the $64$-candidate budget used by comparable
benchmarks, a model used as the mutation operator scores inside the range covered by
random perturbation of a numeric constant, while consuming $1.74$ million output tokens
against the control's none. Held-out performance also degrades after roughly five rounds
of refinement even as the development score continues to improve, so additional iteration
is not merely unhelpful but harmful.

We explain these results structurally. Half of the simple policy space is behaviourally
constant by accident, four constant policies account for $232$ of $460$ family members,
and only $4\%$ of that family beats the best constant. The objective is therefore flat
almost everywhere, which is why hill-climbing and random search both stall and why models
converge on constants. Unlike existing benchmarks in this area, whose problems are drawn
from long-published libraries, our environments are generated from a simulator and cannot
appear in any model's training data.
\end{abstract}

\begin{IEEEkeywords}
benchmarking, code generation, dependability, storage systems, reinforcement learning
\end{IEEEkeywords}

\section{Introduction}

A language model asked to write a controller will write one. It will compile, it will
run, it will contain plausible thresholds and sensible-looking comments, and a reviewer
reading it will usually not be able to say whether it is good. This is the ordinary
condition of deploying generated code in a control path, and the benchmarks available
today do not address it. Pass-rate benchmarks establish that code runs. Benchmarks that
score solution quality compare against the best solution anybody has found, which answers
a different and weaker question.

We take a problem where the best attainable behaviour is not a matter of opinion. During
a storage rebuild, a controller chooses how much device bandwidth to give recovery, every
$300$ seconds, trading rebuild speed against user-visible latency and against the risk of
a second failure while redundancy is reduced. Discretised to four rates and driven by a
fixed workload trace, the problem has a state space small enough that backward induction
returns the exact optimal policy in under $80$ seconds per environment. A generated policy
can therefore be placed on an absolute scale: not ``better than the baseline we had'' but
``captures this fraction of what was attainable.''

Our contributions are:

\begin{itemize}
\item A control benchmark whose reference is an exact optimum computed by dynamic
programming, rather than a best-known or expert solution, over eight environments that
vary workload population, cost structure and device class.
\item Environments that are procedurally generated rather than drawn from a published
problem library, so contamination of the evaluation set is excluded by construction
rather than argued about. \TODO{procedural generation is implemented for the parameter
family but the paper currently reports the eight fixed instantiations}
\item A measurement of iterative refinement against two controls that the existing
literature does not run: the same search loop with the model replaced by a random
mutation, and the same loop with the model's memory removed. Both are necessary to
separate reasoning from sampling, and both change the conclusion.
\item A structural account of why the task defeats both models and blind search, in terms
of the measure of the policy space that is behaviourally constant.
\end{itemize}

\section{Related Work}

\textbf{Algorithm discovery with language models.} FunSearch and its successors
established that a model used as a mutation operator inside an evolutionary loop can
discover programs that improve on human-designed ones. The evidence is strong but it is
evidence about the loop, not about the model in isolation; the number of evaluations is
large and the selection pressure does substantial work.

\textbf{Benchmarks for generated algorithms.} CO-Bench~\cite{cobench} evaluates language
model agents on $36$ combinatorial optimisation problems drawn largely from the
OR-Library, scoring against best-known literature solutions, with $15$ models and nine
agentic frameworks at $64$ iterations. HeuriGym~\cite{heurigym} evaluates heuristics
written by nine models across systems, logistics and biology, feeding execution logs and
costs back into the prompt for iterative refinement, and reports a Quality-Yield Index
normalised so that an expert baseline scores $1$.

Our benchmark differs from both in three respects, and we state the resemblances plainly
because they are close. The scoring construction---normalise so that a reference scores
$1$---is the same idea as HeuriGym's QYI. The iterative-refinement protocol is the same
idea as HeuriGym's feedback loop. What differs is, first, that our reference is provably
optimal rather than best-known or expert-produced, which is possible only because we chose
a problem small enough to solve exactly; second, that our task is sequential
decision-making under partial observability rather than one-shot construction of a
solution to a static instance; and third, that our evaluation varies the environment
rather than the instance, so a policy is tested under cost functions and device
characteristics it was not developed against.

A fourth difference is not about design but about provenance. The OR-Library problems
have been public for decades, together with their best-known solutions and the standard
heuristics for them. A model that produces a strong algorithm for one of them may be
designing or may be recalling, and the benchmark cannot distinguish the two. Our
environments are generated by a simulator written for this work; no solution to them has
ever been published.

\section{Benchmark Design}

\subsection{The control problem}

A disk has failed and the array is rebuilding onto a spare. Every $300$~s the controller
observes a twelve-element vector---foreground read and write load, bandwidth share,
current and target latency, a failure-hazard estimate and its uncertainty, remaining
redundancy, rebuild progress, work remaining, the currently applied rate, and available
headroom---and returns one of four rebuild rates, $\{20, 50, 80, 120\}$~MB/s. Every field
is normalised to approximately $[0,1]$; the controller is told this and is told nothing
about future load. An episode ends when the rebuild completes, a second disk fails, or
$2400$ intervals elapse.

Per-interval cost is
\begin{equation}
c_t = \alpha \, d_t + \beta \, u_t + \gamma \, v_t ,
\end{equation}
where $d_t$ is latency degradation, $u_t$ indicates the array is still degraded, and $v_t$
indicates mean latency exceeded its target. The default weights are
$(\alpha,\beta,\gamma) = (1.0, 0.1, 10.0)$.

Two filters sit between the controller's request and what is applied. A reliability filter
raises the action when redundancy exposure would be too high; a latency filter lowers it
to keep utilisation under the violation threshold. The applied action is therefore
$\max(\text{floor}, \min(\text{requested}, \text{ceiling}))$. Both filters act identically
on every policy including the references, and Section~\ref{sec:saturation} shows that in
one environment they dominate the outcome entirely.

\subsection{Workload}

Foreground load is taken from Tencent Cloud volume traces: $44$ volumes over a shared
$14$-day window, aggregated at $300$~s resolution, selected for the strongest diurnal
cycle. The selection is deliberate and is a limitation we state rather than hide---these
are extremes, not a representative sample of the fleet. The device and failure response
are modelled; in the configuration used here the failure hazard is an unfitted heuristic
whose own documentation forbids describing it as a risk estimate, chosen so that the
benchmark runs from a clean clone with no large artifacts to download.
\TODO{report the fitted-hazard environment as an additional condition, or state its
absence as a limitation}

\subsection{Reference policies and the score}

For each scenario we compute three quantities: $J_{\text{rule}}$, the return of a tuned
risk-aware rule-based controller; $J^\star$, the exact optimum by backward induction over
(rebuild progress $\times$ time); and $J_\pi$, the return of the policy under test. The
score is
\begin{equation}
\G = \frac{J_\pi - J_{\text{rule}}}{J^\star - J_{\text{rule}}} ,
\end{equation}
so $\G = 1$ matches the optimum, $\G = 0$ matches the rule, and negative values are worse
than the rule. $\G$ is dimensionless and invariant to the reward weights, which matters
because those weights are a modelling choice rather than a measurement. A scenario is
excluded when the headroom $J^\star - J_{\text{rule}}$ falls below $0.1\%$ of
$|J_{\text{rule}}|$, since no policy can meaningfully differ there.

\subsection{Environments}

Eight held-out environments vary three axes: workload population (three populations
disjoint from the development set), cost structure (a variant with the violation weight
reduced to $1$, and one with the degraded-state weight raised to $1.0$), and device class
(an SSD-like latency model with a tight target, and an HDD-like one). Mean headroom
differs by a factor of thirty across them, from $2.96$ to $89.04$, which has consequences
we return to in Section~\ref{sec:saturation}.

\subsection{Execution}

Submitted code runs in an isolated subprocess with an import allowlist, no file or network
access, a per-decision time budget and a memory cap. Failures are recorded rather than
retried. A submission that fails to produce a policy function is a recorded zero, not a
missing row. Failures of the harness itself---a sandbox worker that never starts---are
recorded separately and never counted against a model, because on a benchmark whose
headline is a count of working policies, attributing infrastructure failures to models
inflates precisely the statistic being reported.

\section{Baselines}
\label{sec:baselines}

A benchmark of this kind is only as meaningful as its floor. We report four.

\textbf{Uniform random}, selecting an action uniformly at each decision, scores a median
$\G$ of $-99.37$ across the eight environments.

\textbf{Constant policies}, returning one fixed rate and reading no observation at all.
The best of the four, $50$~MB/s, scores $-1.246$.

\textbf{The single-threshold family}: a constant, or one threshold on one observation
field selecting between two actions. We enumerate all $460$ members. The best scores
$-0.377$, which still does not match the tuned rule.

\textbf{A four-line policy} used in early versions of this work as a nonsense floor,
$\texttt{1 if } \sin(\texttt{obs[2]}) > 0.3 \texttt{ else 2}$. It is not nonsense.
Over the normalised range of its input, $\sin$ is monotonic and crosses $0.3$ at $0.305$,
so the policy is exactly a load threshold at $0.305$, and it scores $-0.625$, within noise
of the corresponding explicit threshold at $-0.641$. We report it because the distinction
between ``worse than nonsense'' and ``worse than a one-line rule'' is the difference
between a claim that is false and one that is defensible.

\section{Results}

\subsection{Zero-shot}

Table~\ref{tab:zeroshot} reports the hosted frontier models, five independently sampled
policies each, at a generation budget of $32768$ tokens. We summarise by the median over
the eight environments rather than the mean, because the environments differ in headroom
by a factor of thirty and a mean is decided almost entirely by the environment with the
least room for a policy to matter. \emph{beat\_sine} counts the environments on which a
model exceeds the single-threshold baseline of Section~\ref{sec:baselines}.

\begin{table}[t]
\caption{Zero-shot, hosted frontier models at a $32768$-token generation budget. Median
$\G$ over the eight held-out environments; \emph{valid} counts samples that produced a
runnable policy.}
\label{tab:zeroshot}
\centering
\begin{tabular}{lrrr}
\toprule
Model & valid & median \G & beat\_sine \\
\midrule
\textit{single-threshold baseline} & 1/1 & $-0.62$ & --- \\
\textit{greedy template}           & 1/1 & $-1.68$ & 3/8 \\
\midrule
GPT-6.2-sol                  & 5/5 & $-2.62$  & 1/8 \\
Gemini-3.1-pro-preview       & 5/5 & $-2.82$  & 0/8 \\
DeepSeek-V4.1-flash          & 5/5 & $-9.19$  & 0/8 \\
DeepSeek-V4-Pro              & 5/5 & $-17.54$ & 0/8 \\
Gemini-3.1-flash-lite        & 5/5 & $-28.60$ & 0/8 \\
\bottomrule
\end{tabular}
\end{table}

Every model is below the single-threshold baseline, and all but one are below the
constant-rate policy at $-1.246$. The best model exceeds the baseline on one environment
of eight; four models exceed it on none.

\subsection{A note on generation budget}
\label{sec:budget}

The generation budget is part of the protocol, not an implementation detail: a model that
exhausts it is recorded as producing no policy, which is a property of the constant rather
than of the model. We therefore report only models that shared a budget. An earlier
open-weight roster was generated at $8192$ tokens and is excluded from
Table~\ref{tab:zeroshot}. The exclusion is not precautionary. Inspecting the retained raw
replies, DeepSeek-R1-Distill-32B's terminate mid-expression---one ends
\texttt{if work\_remaining >0.}---so the budget bound for that model, and one of its five
samples produced no policy for that reason alone. Because every reply is retained, this is
auditable rather than inferred. \TODO{regenerate the open-weight roster at 32768 and
report it as a second table; blocked on cluster access}

\subsection{Iterative refinement}

We ran a conversational loop for ten rounds: the model writes a policy, the policy is
scored on the development workload, and the measurements are returned as the next turn.
The eight evaluation environments are scored every round and recorded but never shown.
Two arms differ in one respect: whether the conversation is carried forward
(\texttt{full}) or discarded between rounds (\texttt{none}).

Table~\ref{tab:iterative} reports, for each arm, the held-out score of whichever round
scored best on development---the selection rule a practitioner would actually use.

\begin{table}[t]
\caption{Iterative refinement, best-on-development selection. Held-out $\G$ is the median
over the eight environments. Random search is over the single-threshold family at a
matched budget of evaluations.}
\label{tab:iterative}
\centering
\begin{tabular}{lr}
\toprule
Method, 10 attempts & Held-out \G \\
\midrule
Random search, 10 evaluations      & $-1.445$ \\
DeepSeek-V4-Pro, full memory       & $-1.43$ \\
DeepSeek-V4-Pro, no memory         & $-0.43$ \\
GPT-6.2-sol, no memory\footnotemark[1] & $-0.64$ \\
GPT-6-luna, full memory            & $-14.63$ \\
GPT-6-luna, no memory              & $-17.47$ \\
\bottomrule
\end{tabular}
\footnotetext[1]{Two repeats at a different generation budget; not matched.}
\end{table}

Three observations. Carrying the conversation produces $-1.43$, and blind random search at
the same budget produces $-1.445$: indistinguishable. Deleting the model's memory of its
own earlier attempts \emph{improves} the result by a factor of three. And taking the final
round rather than the best round is catastrophic---$-11.78$ against $-1.58$ with memory,
$-52.58$ against $-0.69$ without---so the gain that iteration appears to deliver is
produced by selection, not by refinement.

Held-out performance degrades with further iteration. In four of the five conversations
where we measured rounds $1$, $5$ and $10$, round $10$ generalises worse than round $5$;
in one case the held-out score moves from $-0.28$ to $-171.04$ while the development score
the model was optimising continued to look acceptable.

\subsection{Population search}

Because the conversational loop is not what the comparable benchmarks run, we repeated the
comparison inside a population search: keep a pool of the best candidates, draw a parent,
produce a child, score it, insert. Only the operator producing the child differs---a
language model asked to improve the parent, a random perturbation of one numeric constant
in the parent's own source, or a fresh draw from the simple family. Budget is $64$ scored
candidates, matching CO-Bench.

\begin{table}[t]
\caption{Population search at a budget of $64$ scored candidates, matching CO-Bench.
Development score of the best candidate found, the held-out score of that same candidate,
and the output tokens the arm consumed.}
\label{tab:evolve}
\centering
\begin{tabular}{lrrrr}
\toprule
Arm & dev \G & held-out \G & valid & output tokens \\
\midrule
mutate, seed 0 & $-1.588$ & $-1.246$ & 64/64 & $0$ \\
mutate, seed 1 & $-1.588$ & $-1.246$ & 64/64 & $0$ \\
mutate, seed 2 & $-1.588$ & $-1.384$ & 64/64 & $0$ \\
sample, seed 0 & $-1.588$ & $-1.246$ & 64/64 & $0$ \\
sample, seed 1 & $-1.588$ & $-1.548$ & 64/64 & $0$ \\
sample, seed 2 & $\mathbf{-0.390}$ & $-0.708$ & 63/64 & $0$ \\
\midrule
DeepSeek-V4.1-flash, seed 0 & $-0.553$ & $-0.447$ & 52/64 & $1{,}744{,}920$ \\
\bottomrule
\end{tabular}
\end{table}

Table~\ref{tab:evolve} gives the result. Both control operators converge to $-1.588$ on
development---the best constant policy---on five of six seeds, typically by the eighth
candidate. The exception is informative: one \texttt{sample} seed found $-0.390$, which is
the exact optimum of the single-threshold family, so blind sampling occasionally solves
the restricted space outright.

The language model's $-0.553$ falls \emph{inside} the control range of $-1.588$ to
$-0.390$, and is worse than the best control seed. With one language-model seed against
three seeds per control operator we do not claim a difference in either direction; we
claim the absence of a detectable one. \TODO{two further seeds of the reasoning model and
three of a non-reasoning model are running; update this paragraph and the table}

What does separate is cost. The model consumed $1.74$ million output tokens and
approximately four hours of wall-clock to reach a score the controls reach in ten minutes
with no model at all, and it produced unrunnable code on $12$ of $64$ candidates against
roughly one for the controls. It was also behind for most of the run, holding at $-3.21$
through candidate $32$ and only reaching $-0.553$ at candidate $60$; truncating the budget
slightly below CO-Bench's would have recorded it as clearly worse.

\section{Why the Task Is Hard}
\label{sec:saturation}

The control results are not a property of weak baselines. They follow from the shape of
the objective.

Of the $460$ members of the single-threshold family, $232$---exactly half---score one of
four values, and those four values are the scores of the four constant policies. The
reason is that a threshold placed where the observation never crosses it yields a policy
that is behaviourally constant: \texttt{3 if obs[2] > 0.92 else 1} is simply ``always
$1$''. Only $19$ members, $4\%$, beat the best constant, and they occupy narrow threshold
bands. A hill-climber inside a plateau receives an identical score for every perturbation
and has nothing to climb, which is why both control operators stall at the best constant
and why models converge on constant policies rather than because they are lazy.

A second mechanism operates in one environment. The reliability floor and the latency
ceiling compose as $\max(\text{floor}, \min(\text{requested}, \text{ceiling}))$, so when
the ceiling is pinned low almost everywhere, every request above it collapses to the same
applied action. In the SSD-class environment, whose mean headroom of $2.96$ is four times
smaller than the next smallest, two submissions with different action distributions and
different per-decision timings produce violation fractions identical to sixteen
significant figures and bit-identical per-scenario returns: the filters erased every one of
one policy's $416$ deviations before they reached the trajectory. That environment measures
the safety architecture, not the policy, and we report it as a stress probe rather than as
a ninth ordinary environment. \TODO{report the clipped-decision fraction per environment
once the instrumentation lands}

\section{Threats to Validity}

\textbf{One simulator.} The eight environments are eight parameterisations of a single
simulator. They vary workload, cost structure and device response, but they do not vary
the underlying dynamics, and a result that holds across them may not transfer to a
different recovery mechanism.

\textbf{Modelled failure process.} The hazard model in the configuration reported here is
an unfitted heuristic. A fitted model exists in this project and is not wired into the
benchmark, by a deliberate choice to keep the artifact small.

\textbf{Selected workload.} The $44$ volumes were chosen for the strongest diurnal cycle
and are extremes rather than a representative sample.

\textbf{Iteration budget.} We report ten rounds of conversational refinement and $64$
candidates of population search. CO-Bench reports gains at $64$ iterations across nine
frameworks; our negative result is scoped to the loops we ran and does not establish that
no agentic framework can do better.

\textbf{Scale of the model evidence.} Three repeats per arm, two language models in the
iterative track. The variance between repeats is large and the conclusions about memory
rest on a narrow base. \TODO{additional models if budget permits}

\textbf{Hosted inference.} Models were served through a routing provider. An earlier batch
in which repeats were served by different hosts produced bimodal results that vanished
once each model was pinned to a single provider, so we pin, and we note that scores from
hosted endpoints are not reproducible in the way local weights are.

\section{Conclusion}

Given a task where the optimum is computable, generated control code can be measured
rather than reviewed, and what the measurement shows is uncomfortable: policies that run
and read sensibly can be beaten by a constant; refinement loops deliver what sampling
delivers; and the obvious selection rule, taking the last iteration, is far worse than
keeping the best. None of this is visible from pass rates, and none of it would have been
visible without a reference that is optimal rather than merely good.

\begin{thebibliography}{9}
\bibitem{cobench} W.~Sun, S.~Feng, S.~Li, and Y.~Yang, ``CO-Bench: Benchmarking Language
Model Agents in Algorithm Search for Combinatorial Optimization,'' AAAI, 2026.
\bibitem{heurigym} H.~Chen \emph{et al.}, ``HeuriGym: An Agentic Benchmark for LLM-Crafted
Heuristics in Combinatorial Optimization,'' ICLR, 2026.
\TODO{FunSearch, AlphaEvolve, Eureka, OR-Library, Tencent trace, Backblaze}
\end{thebibliography}

\end{document}
